% Tema de la presentación. Toma los colores de la propia aplicación, que es
% oscura, para que las capturas no queden como recuadros negros flotando sobre
% blanco y el conjunto se lea como una sola cosa.

\usetheme{default}
\usecolortheme{default}
\setbeamertemplate{navigation symbols}{}

% Paleta copiada de frontend/src/index.css y de la de grupos de App.jsx.
\definecolor{fondo}{HTML}{0B0C0E}
\definecolor{superficie}{HTML}{14171B}
\definecolor{superficiedos}{HTML}{1B1F25}
\definecolor{linea}{HTML}{2A2F36}
\definecolor{texto}{HTML}{E8E9EB}
\definecolor{apagado}{HTML}{8E949D}
\definecolor{tenue}{HTML}{565C64}
\definecolor{acento}{HTML}{4C7DF0}
\definecolor{acentoclaro}{HTML}{6F97F5}
\definecolor{verde}{HTML}{58C4A3}
\definecolor{ambar}{HTML}{E0A458}
\definecolor{rojo}{HTML}{D6685E}
\definecolor{morado}{HTML}{A98BD0}

\setbeamercolor{background canvas}{bg=fondo}
\setbeamercolor{normal text}{fg=texto, bg=fondo}
\setbeamercolor{structure}{fg=acento}
\setbeamercolor{alerted text}{fg=acentoclaro}
\setbeamercolor{frametitle}{fg=texto, bg=fondo}
\setbeamercolor{itemize item}{fg=acento}
\setbeamercolor{enumerate item}{fg=acento}

\setbeamerfont{frametitle}{size=\Large, series=\bfseries}
\setbeamerfont{footline}{size=\tiny}

\setbeamertemplate{frametitle}{%
	\vspace{0.5em}\hspace{0.055\paperwidth}%
	\begin{minipage}{0.89\paperwidth}
		\raggedright\usebeamerfont{frametitle}\color{texto}%
		\insertframetitle\par\vspace{0.22em}%
		{\color{acento}\rule{2.6em}{2.5pt}}
	\end{minipage}\par\vspace{0.25em}}

% El número que se muestra es el de página del PDF, que es el que aparece en el
% visor y el que usa el guion. Así no hay dos numeraciones distintas.
\setbeamertemplate{footline}{%
	\hbox{\begin{beamercolorbox}[wd=\paperwidth, ht=1.9ex, dp=1.2ex,
			leftskip=0.055\paperwidth, rightskip=0.055\paperwidth]{}%
		\usebeamerfont{footline}\color{tenue}%
		\hfill\thepage%
	\end{beamercolorbox}}}

\setbeamertemplate{itemize item}{\color{acento}\raisebox{0.15ex}{\rule{0.42em}{0.42em}}}
\setbeamertemplate{itemize subitem}{\color{tenue}\raisebox{0.18ex}{\rule{0.3em}{0.3em}}}
\setbeamertemplate{enumerate item}{\color{acento}\bfseries\insertenumlabel}
\setlength{\leftmargini}{1.3em}

% ---------------------------------------------------------------------------
% Separador de bloque: el número muy grande, en color, y el nombre debajo.
\AtBeginSection[]{%
	\begin{frame}[plain, noframenumbering]
		\vfill
		\hspace{0.09\paperwidth}%
		\begin{minipage}{0.8\paperwidth}
			{\color{acento}\fontsize{72}{72}\selectfont\bfseries\insertsectionnumber}\\[0.25em]
			{\color{texto}\fontsize{26}{30}\selectfont\bfseries\insertsectionhead}\\[0.7em]
			{\color{acento}\rule{3.5em}{2.5pt}}
		\end{minipage}
		\vfill
	\end{frame}}

% ---------------------------------------------------------------------------
% Diapositiva con una captura ocupando toda la superficie.
\newenvironment{lienzo}[1]{%
	\begingroup
	\usebackgroundtemplate{\includegraphics[width=\paperwidth, height=\paperheight]{#1}}%
	\begin{frame}[plain]}%
	{\end{frame}\endgroup}

% Caja de texto que se superpone a una captura. El fondo semitransparente deja
% ver la escena por detrás sin que el texto pierda contraste.
\newcommand{\sobreimagen}[3][north west]{%
	\begin{tikzpicture}[remember picture, overlay]
		\node[anchor=#1, fill=fondo, fill opacity=0.9, text opacity=1, draw=linea,
		      rounded corners=6pt, inner sep=1.1em, text width=0.4\paperwidth, align=left,
		      font=\color{texto}] at ([shift={(#2)}]current page.#1) {#3};
	\end{tikzpicture}}

% Título grande sobre una captura, sin caja.
\newcommand{\tituloimagen}[2]{%
	\begin{tikzpicture}[remember picture, overlay]
		\node[anchor=west, inner sep=0pt] at ([shift={(0.075\paperwidth,0)}]current page.west) {%
			\begin{minipage}{0.55\paperwidth}
				\color{texto}#1\\[0.5em]{\color{acento}\rule{3em}{2.5pt}}\\[0.6em]
				\color{apagado}#2
			\end{minipage}};
	\end{tikzpicture}}

% Velo oscuro para que el texto se lea sobre la captura.
\newcommand{\velo}[1][0.72]{%
	\begin{tikzpicture}[remember picture, overlay]
		\fill[fondo, opacity=#1] (current page.south west) rectangle (current page.north east);
	\end{tikzpicture}}

% Velo solo en la mitad izquierda, degradado, para la portada.
\newcommand{\veloizq}{%
	\begin{tikzpicture}[remember picture, overlay]
		\fill[fondo] (current page.south west) rectangle
			([xshift=0.44\paperwidth]current page.north west);
		\fill[fondo, path fading=east]
			([xshift=0.44\paperwidth]current page.south west) rectangle
			([xshift=0.78\paperwidth]current page.north west);
	\end{tikzpicture}}

% ---------------------------------------------------------------------------
% Tarjeta: bloque con fondo propio para agrupar una idea.
\newcommand{\tarjeta}[3][acento]{%
	\begin{tikzpicture}
		\node[fill=superficie, draw=linea, rounded corners=5pt, inner sep=0.7em,
		      text width=\dimexpr\linewidth-1.9em, align=left] {%
			{\color{#1}\bfseries\small #2}\\[0.3em]{\color{apagado}\footnotesize #3}};
	\end{tikzpicture}}

% Variante de una sola línea, para cuando hay muchas tarjetas juntas.
\newcommand{\tarjetafina}[3][acento]{%
	\begin{tikzpicture}
		\node[fill=superficie, draw=linea, rounded corners=5pt, inner sep=0.6em,
		      text width=\dimexpr\linewidth-1.7em, align=left] {%
			{\color{#1}\bfseries\small #2}\ {\color{apagado}\footnotesize #3}};
	\end{tikzpicture}}

% Cifra destacada con su etiqueta debajo.
\newcommand{\cifra}[3][acento]{%
	\begin{minipage}[t]{0.3\textwidth}\centering
		{\color{#1}\fontsize{30}{32}\selectfont\bfseries #2}\\[0.3em]
		{\color{apagado}\small #3}
	\end{minipage}}

% Los diagramas de la memoria tienen fondo blanco, así que se colocan sobre una
% tarjeta clara en lugar de dejarlos recortados sobre el fondo oscuro.
\newcommand{\diagrama}[2][0.66\textheight]{%
	\begin{tikzpicture}
		\node[fill=white, rounded corners=6pt, inner sep=0.6em]
			{\includegraphics[height=#1, keepaspectratio]{#2}};
	\end{tikzpicture}}

% Rótulo sobre una captura a sangre. Poco texto y bien grande, que es lo que
% se lee desde el fondo de un aula.
\newcommand{\rotulo}[1]{%
	\begin{tikzpicture}[remember picture, overlay]
		\node[anchor=north west, fill=fondo, fill opacity=0.88, text opacity=1, draw=linea,
		      rounded corners=8pt, inner sep=0.9em, align=left, font=\color{texto}]
			at ([shift={(0.05\paperwidth,-0.08\paperheight)}]current page.north west) {#1};
	\end{tikzpicture}}

\newcommand{\rotuloder}[1]{%
	\begin{tikzpicture}[remember picture, overlay]
		\node[anchor=north east, fill=fondo, fill opacity=0.88, text opacity=1, draw=linea,
		      rounded corners=8pt, inner sep=0.9em, align=left, font=\color{texto}]
			at ([shift={(-0.05\paperwidth,-0.08\paperheight)}]current page.north east) {#1};
	\end{tikzpicture}}

% Frase suelta, grande y centrada: una idea por diapositiva.
\newcommand{\frase}[2][texto]{%
	\vfill
	\begin{center}
		\begin{minipage}{0.8\paperwidth}\centering
			{\color{#1}\fontsize{22}{28}\selectfont\bfseries #2}
		\end{minipage}
	\end{center}
	\vfill}

\newcommand{\clave}[1]{{\color{acentoclaro}\bfseries #1}}
\newcommand{\suave}[1]{{\color{apagado}#1}}
